\section{Introduction}
\label{sec:intro}

Several quadrotors sharing a payload on cables can lift what none of them could carry alone, and spread the load so
that each flies with margin \cite{olleroPresentFutureAerial2022, estevezReviewAerialTransportation2024}. Such a team
is almost always treated as fixed: the carriers are attached on the ground, take off together and land together. A
team that could change in the air would be more useful and more robust. A carrier with a low battery or a fault could
leave and come back; a fresh carrier could join to take up a heavier load or relieve another; and a carrier that is
lost without warning, a released magnet or a failed motor, need not bring the payload down with it. This letter
asks what a controller needs to let a team carrying a cable-suspended payload lose and gain carriers in flight.

\emph{Cooperative cable transport.} The statics of a rigid payload hung from several cables, and the internal forces
the team may choose freely, are well understood \cite{michael2011cooperative, tognon2018internal,
gabellieri2023equilibria}. Controllers range from decentralised schemes that need no communication between carriers
\cite{tognon2018internal, gabellieri2023force, klausen2020cooperative} through geometric and leader--follower designs
\cite{lee2018geometric, loianno2018cooperative} to the passivity-based dissipative network of Quan et
al.~\cite{quan2025swarm}. Model predictive control (MPC) suits the problem because cable tensions are bounded and
the payload's motion couples every carrier: nonlinear MPC of the full payload pose \cite{sunNonlinearMPCFullPose2023,
liNonlinearModelPredictive2023}, kinodynamic planning for the team \cite{wahbaKinodynamicMotionPlanning2024}, and
distributed formulations \cite{decarliDistributedNMPCCooperative, distributedNMPC} all fly a team of a fixed size.

\emph{Reconfiguration and attach.} Plug-and-play MPC adds and removes subsystems from a distributed controller
\cite{riverso2013plug, zeilinger2013plug, aboudonia2022reconfigurable, saccani2023model}, but for networks of
loosely coupled plants, not for carriers that share one payload. Modular aerial robots dock and undock in flight
\cite{saldana2018modquad, sugihara2023trady, sugihara2026legion}, and magnetic grippers let a multirotor pick up and
release objects \cite{gawel2017magnetic, miyazaki2018airborne}. Fault-tolerant cooperative transport keeps a payload
up after a carrier fails \cite{liang2021fault, farrasFaultTolerantControlMultiQuadcopter2024}. We know of no work
in which a team carrying a cable-suspended payload both loses carriers without warning and admits carriers that join
in the air, nor of one flown on hardware.

The working hypothesis of this letter is that a centralised optimal control problem suits a change of team because
the planner solves for the force distribution of the team it actually has in the first tick after the event,
instead of leaving the survivors to find it from the motion that follows. The contributions are:
\begin{itemize}
  \item a cooperative MPC for a cable-suspended ring whose attachment geometry is a runtime parameter, so that a
  change of team is a switch to a pre-built problem of the new size, warm-started from the measured state
  (Section~\ref{sec:method});
  \item detection of an unannounced release from motion capture alone, a drop rule from the support margin of the
  survivors, and the rejoin of a departed carrier to its own plate in flight (Sections~\ref{sec:reconfig}
  and~\ref{sec:rejoin});
  \item the attach of a newcomer flown by another team's controller, admitted only on a motion-capture proof that its
  magnet moves with the ring;
  \item validation in a digital twin built from the rig's identified thrust map and geometry and on a
  motion-capture rig with unmodified Betaflight quadrotors, with a scalability study and a comparison against our
  port of Quan et al.'s method \cite{quan2025swarm} (Sections~\ref{sec:res_reconfig}--\ref{sec:res_compare}).
\end{itemize}
